\documentclass[Screen16to9,17pt]{foils}
\usepackage{zencurity-slides}

\externaldocument{packet-construction-workshop-exercises}
\selectlanguage{english}
\usepackage{listings}

\begin{document}

\mytitlepage{MoonGen BlueField-2 Packet Generation}{BornHack July 2026}

%\hlkprofiluk

\slide{Experiment Overview}
Goal: reproduce and extend the 2025 version of the DPDK/MoonGen packet generation results on hosts \verb+T150+ and \verb+basse+ equipped with an identical BlueField-2 DPU network cards

Both hosts use the \verb+penguinping+ MoonGen scripts against a TCP SYN target, transmitting via the BF2's integrated ConnectX-6 Dx controller

Source for the tools:
\begin{list2}
\item \url{https://github.com/tumi8/MoonGen.git} TUM School of Computation, Information and Technology
Technical University of Munich
\item \url{https://codeberg.org/kramse/penguinping} My Lua scripting to simulate hping3 using Lua
\item \url{https://penguinping.org} Home page for the tool Lua scripting to simulate hping3 using Lua -- to be updated!!!!111
\end{list2}


\slide{Penguinping testing}

\hlkimage{16cm}{ddos-test-layout.png}

\slide{Penguinping Software stack}

\begin{itemize}
\item I have previously used a nice library Libmoon and MoonGen for creating packets with Lua and 10G cards -- one CPU core could push 14million packets per second
\item Original paper: P.~Emmerich et al., \emph{MoonGen: A Scriptable High-Speed Packet Generator}, ACM IMC 2015. \url{https://dl.acm.org/doi/10.1145/2815675.2815692}
\item Updated fork with MoonEm path emulator (delay, rate limiting, packet loss): TU Munich Network Architectures and Services group, CoNEXT 2025. \url{https://github.com/tumi8/moongen}
\item libmoon (core DPDK/LuaJIT framework underlying MoonGen): \url{https://github.com/libmoon/libmoon}
\item This is confirmed working on Debian 13 in 2026!
\item It used DPDK, Just-in-time compilation and is very nice
\end{itemize}

I promise to update \url{https://penguinping.org/} and \url{https://codeberg.org/kramse/penguinping}


\slide{PenguinPing -- re-implementing hping3 with Lua}

\begin{alltt}\footnotesize
root@penguin01:~/projects/MoonGen# {\bfseries ./build/MoonGen ./examples/penguinping-02.lua 10.0.49.1 -a 10.1.2.3 -r 10000 -S -p 80}
[INFO]  Initializing DPDK. This will take a few seconds...
[INFO]  Found 2 usable devices:
   Device 0: 00:25:90:32:9F:F2 (Intel Corporation 82599ES 10-Gigabit SFI/SFP+ Network Connection)
[INFO]  Device 0 (00:25:90:32:9F:F2) is up: 10000 MBit/s
TCP mode get TCP packet
IP4 10.1.2.3 > 10.0.49.1 ver 4 ihl 5 tos 0 len 46 id 0 flags 0 frag 0 ttl 64 proto 0x06 (TCP) cksum 0x0000 [-]
TCP 52049 > 80 seq# 1 ack# 0 offset 0x5 reserved 0x00 flags 0x02 [-|-|-|-|SYN|-] win 10 cksum 0x0000 urg 0 []
  0x0000:   0000 0000 0000 0025 9032 9ff2 0800 4500
  0x0010:   002e 0000 0000 4006 0000 0a01 0203 0a00
  0x0020:   3101 cb51 0050 0000 0001 0000 0000 5002
  0x0030:   000a 0000 0000 0000 0000 0000

[Device: id=0] TX: 14.88 Mpps, 7619 Mbit/s (9999 Mbit/s with framing)
[Device: id=0] TX: 14.78 Mpps, 7568 Mbit/s (9933 Mbit/s with framing)
[Device: id=0] TX: 14.88 Mpps, 7619 Mbit/s (10000 Mbit/s with framing)
\end{alltt}

\begin{list2}
\item Using Lua we can implement the same attacks from Hping3 easily
\item Only about 230 lines of Lua using MoonGen and libmoon
\item Can run at specific rate up to full 10Gbps / 14.8 Million packets per second using a single CPU core
\end{list2}


\slide{Comparable to real DDoS?}

Tools are simple and widely available but are they actually producing same result
as high-powered and advanced criminal botnets. We can confirm that the attack
delivered in this test is, in fact, producing the traffic patterns very close to
criminal attacks in real-life scenarios.

\begin{list2}
\item We can also monitor logs when running a single test-case
\item Gain knowledge about supporting infrastructure
\item Can your syslog infrastructure handle 800.000 events in $<$ 1 hour?
\end{list2}

Main difference are that attackers are free to switch attack types and mix them. While we try specifically to keep using one type, to see the worst and which ones that hurt the most.

I also start at the bottom, and work my way up -- while an attacker may begin attacking HTTP/HTTPS directly.


\slide{Nvidia Mellanox DPU2}

\hlkimage{12cm}{nvidia-dpu2.jpg}

\begin{itemize}
\item MT42822 — ConnectX-6 Dx network controller PCIe Gen4 host interface
\item 2× 100G ports (SFP28/DAC)
\item 8× Arm Cortex-A72 onboard cores -- not used in this experiments, slower than main CPU
\item Drivers: \verb+mlx5_core+ (host), \verb+mlx5+ PMD (DPDK)
 \end{itemize}


\slide{Hardware: Small Dell PC architecture}

\begin{list2}
\item Dell PowerEdge T150: Intel Xeon E-2314, 4 cores (AVX-512 capable)\\
Debian 13 (trixie), kernel 6.12.90+deb13.1-amd64
BlueField-2 DPU, MT42822, integrated ConnectX-6 Dx, dual 100G ports\\
\item Dell Precision 3640 Basse: Intel Core i9-10900K, 10 physical cores / 20 threads (no AVX-512)\\
Debian 13 (trixie), kernel 6.12.95+deb13-amd64
BlueField-2 DPU, same model as T150\\
\item Port 0 on both systems connected to a Mellanox SN2700 switch
\end{list2}

\slide{Software Stack}
\begin{list2}
\item NVIDIA DOCA-Host 3.4.0 (\verb+doca-host_3.4.0-085000-26.04-debian13_amd64.deb+) on both hosts
\item DPDK via \verb+dpdk-community+ (DOCA-bundled), MoonGen built from a fresh clone against each host's vendored DPDK
\item MLNX-OFED DKMS modules (\verb+mlnx-ofed-kernel+, \verb+iser+, \verb+isert+, \verb+srp+, \verb+xpmem+) installed and matching between hosts
\item A bit of fiddling with kernel options \verb+/etc/default/grub/+:\\
\verb+GRUB_CMDLINE_LINUX_DEFAULT="quiet intel_iommu=off default_hugepagesz=2MB hugepagesz=2MB hugepages=2048"+
\end{list2}


\slide{Initial Results}

T150 Results: Verified\\
\begin{tabularx}{\textwidth}{lX}
Queues & 3 TX queues, 3 of 4 CPU cores \\
Peak throughput & approximately 43 Mpps \\
Line rate & approximately 20 Gbps \\
Frame size & 60 bytes \\
Rate limiting & hardware \verb+setRate+ API unsupported on BF2, silently ignored \\
\end{tabularx}

Basse Results: Verified\\
\begin{tabularx}{\textwidth}{lX}
Queues & 9 TX queues, 9 of 10 physical cores \\
Peak throughput & approximately 97 to 98 Mpps sustained \\
Verification & hardware phy counters (host) and switch port counters (SN2700) both confirm real traffic \\
\end{tabularx}

Combined traffic would be close to the maximum for a 100Gbps port with small packets ~140Mpps


\slide{Packet Generation Summary — MoonGen on BlueField-2}
\begin{itemize}
\item Hosts: Dell PowerEdge T150, Intel Xeon E-2314 (4 cores, 2.80 GHz)\\
 Dell Precision 3640
Intel Core i9-10900K, 10 physical cores / 20 threads (no AVX-512)
\item NIC: NVIDIA BlueField-2, single 100G port
\item Software: MoonGen + libmoon updated by TUM School of Computation, Information and Technology
Technical University of Munich
, DPDK \verb+mlx5+ PMD
\item Packet size: 60 bytes (minimum Ethernet frame)
\item TX queues: 3 (1 core reserved for DPDK master)
\item \textbf{Result: \textasciitilde{}43 Mpps / \textasciitilde{}20 Gbps}
\item Verified via \verb+ethtool -S+ hardware counters (\verb+rx_packets_phy+)
\item Bottleneck: CPU cores — scales linearly with queue count
\end{itemize}

It's in my tent on PTY, if you want to try it! aaaand I brought a 10-core CPU PC



\slide{Context: ConnectX-6 Dx Ceiling}
Independent DPDK community benchmarks on the same ConnectX-6 Dx silicon report a single-port ceiling of approximately 148 Mpps at 64-byte frames, and approximately 210 to 214 Mpps combined when both ports of a dual-port card are driven simultaneously

Basse's verified approximately 98 Mpps single-port result sits at roughly two thirds of that single-port ceiling, with core count (9 of 10 physical cores already committed to TX) as the likely next constraint rather than the NIC itself

Next Steps:\\
\begin{list2}
\item Push queue count and burst/cache tuning further on \verb+basse+ to probe the single-port ceiling
\item Connect basse's second BF2 port for genuine dual-port testing, budgeting CPU cores accordingly
\item Make more Lua scripting to expand scenarios for testing -- what about 100Mpps of QUIC?
\end{list2}

\myquestionspage

\end{document}
